\documentclass[10pt,a4paper]{amsart}
\usepackage[margin=19mm]{geometry}
\usepackage{amsmath,amssymb,mathtools}
\usepackage[scr=rsfs,cal=euler]{mathalfa}
\usepackage{xfrac}
\usepackage[unicode,hidelinks,bookmarks=false]{hyperref}
\newcommand{\ie}{i.e.\ }
\newcommand{\cf}{cf.\ }
\newcommand{\resp}{respectively\ }
\newcommand{\Subsection}{\subsection}
\providecommand{\nobreakdash}{\nobreak\hbox{-}}
\setlength{\emergencystretch}{2em}
\multlinegap0pt
\usepackage{bm}
\usepackage{bbm}
\usepackage{dsfont}
\usepackage{xspace}
\usepackage{cancel}
\usepackage{mathtools}
\usepackage{amsthm}
\makeatletter
\@ifundefined{assumption}{\newtheorem{assumption}{Assumption}}{}
\@ifundefined{lemma}{\newtheorem{lemma}{Lemma}}{}
\makeatother
\title[Penalty-scaling effects in nonsymmetric interior-penalty DG ]{Penalty-scaling effects in nonsymmetric interior-penalty DG discretizations of viscous rotating shallow-water equations}
\author{Xue Zhang, Jingmin Xia, Xu Qian}
\date{Source: arXiv:2607.15643v1 (CC BY 4.0)}
\begin{document}
\maketitle

\noindent\emph{Source and licence.} Penalty-scaling effects in nonsymmetric interior-penalty DG discretizations of viscous rotating shallow-water equations. Version arXiv:2607.15643v1 (CC BY 4.0), distributed under the Creative Commons Attribution 4.0 International licence, as recorded in the arXiv licence metadata for this exact version and shown on the abstract page https://arxiv.org/abs/2607.15643v1. Full rights notice: licenses/arxiv-2607.15643-CC-BY-4.0.txt.\par\smallskip

\noindent\textit{Editorial scope.} Counted unit: the Lemma whose source label is \texttt{lma:coercivity} in the article (source0.tex lines 582--597, source label \texttt{lma:coercivity}), delivered together with the article's own complete original proof of it (source0.tex lines 599--672, one \texttt{proof} environment of 74 lines). No mathematics is written, repaired or re-derived: statement and proof are the article's own text line for line. Required context retained uncounted: eq7 (lines 337--346); asm:regularity (lines 472--479). Every definition, display and label of those lines is retained unchanged. Omitted from this package and not redistributed: 1--336, 347--471, 480--581, 598--598, 673--1559. Nothing inside a retained excerpt is omitted or altered: every retained body is delivered exactly as the article's lines are written, so no omission receipt of any kind accompanies this package. Citations: the retained excerpts carry no citation command at all, so no bibliography entry of the article is transcribed and no reference is redistributed. Reference adaptation: none was needed and none was made; every reference inside the delivered unit resolves inside it, so no unresolved reference or citation is produced. Printed slips kept literal: no printed word, symbol, hypothesis or number of the counted statement or of its proof was altered. Layout and class: the article's own class is \texttt{sn-jnl} with packages including fontenc, inputenc, amsmath, amssymb, amsfonts, amsthm, mathtools, bm, graphicx, subcaption, booktabs, multirow, enumitem, xcolor; the wrapper is the lane's pinned \texttt{amsart} editorial wrapper at 10pt on A4, which restates the theorem-like environments the retained text needs and adds the article's own macro definitions that the retained text needs (none), with extra wrapper packages (bm, bbm, dsfont, xspace, cancel, mathtools). The delivered numbering is the unit's own. Assets: no retained span contains a figure environment, a graphics-inclusion command, a PDF-page inclusion, a table or a TikZ picture; no image file is copied into this package, the package declares no asset dependency and no image file is redistributed with it. Licence: version arXiv:2607.15643v1 of this article is distributed under the Creative Commons Attribution 4.0 International licence, as recorded in the arXiv licence metadata for this exact version and shown on the abstract page https://arxiv.org/abs/2607.15643v1. A full rights notice is delivered at licenses/arxiv-2607.15643-CC-BY-4.0.txt. Characters: every retained excerpt is pure ASCII, so the delivered engine, XeLaTeX with its default Latin Modern text font, reports no missing character.
\begin{equation}
\label{eq7}
\begin{aligned}
a_h(\mathbf z,\mathbf w)
={}&\sum_{K\in Q_h}\int_K\nabla S(\mathbf z):\nabla\mathbf w\,\mathrm d\mathbf x
-\sum_{e\in\mathcal E_h}\int_e\{\!\{\nabla S(\mathbf z)\}\!\}\mathbf n\cdot[[\mathbf w]]\,\mathrm ds\\
&+\sum_{e\in\mathcal E_h}\int_e\{\!\{\nabla\mathbf w\}\!\}\mathbf n\cdot[[S(\mathbf z)]]\,\mathrm ds
+\sum_{e\in\mathcal E_h}\int_e\mu_e[[S(\mathbf z)]]\cdot[[\mathbf w]]\,\mathrm ds,
\end{aligned}
\end{equation}

\begin{assumption}[Regularity and mesh]
\label{asm:regularity}
The exact solution satisfies $\mathbf{q}\in \left(H^{k+1}(\Omega)\right)^3$ for some $k\ge 1$. The mesh family $Q_h$ is shape-regular and quasi-uniform. The reference depth remains strictly positive,
\[
H_0:=\inf_{\Omega}(\eta_0+b)>0.
\]
The penalty prefactor satisfies $\sigma > 0$.
\end{assumption}

\begin{lemma}[Continuity and coercivity]
\label{lma:coercivity}
Under Assumption~\ref{asm:regularity}, and for $\beta\ge 1$, there exists a constant $C_b=C \nu_T \cdot \max \left\{1,\frac{C_k \sqrt{n_0}}{\sqrt{\sigma}}\right\}>0$ independent of $h$ such that
\begin{equation}
\label{eq13}
|a_h(\mathbf{z}_h,\mathbf{w}_h)|
\le C_b\|\mathbf{z}_h\|_{\mathrm{DG}}\|\mathbf{w}_h\|_{\mathrm{DG}},
\quad
\forall \mathbf{z}_h,\mathbf{w}_h\in\mathbf{V}_h.
\end{equation}
For any $\sigma>0$, the coercivity property is unconditionally satisfied
\begin{equation}
\label{eq14}
a_h(\mathbf{z}_h,\mathbf{z}_h)=\nu_T\|\mathbf{z}_h\|_{\mathrm{DG}}^2.
\end{equation}
\end{lemma}

\begin{proof}
We first prove continuity. By the definition of the NIPG bilinear form $a_h$ in \eqref{eq7}, the viscous contribution consists of a volume gradient term and two interface terms, together with the penalty stabilization term. For $\mathbf{z}_h=(\zeta_h,m_{1,h},m_{2,h}),\mathbf{w}_h=(\xi_{h},w_{1,h},w_{2,h}) \in \mathbf V_h$, we write
\begin{equation}\notag
\begin{aligned}
a_h(\mathbf{z}_h,\mathbf{w}_h) \triangleq
I_1+I_2+I_3+I_4,
\end{aligned}
\end{equation}
where
\begin{equation}\notag
\begin{aligned}
I_1 &= \sum_{K \in Q_h}\int_K \nabla S(\mathbf{z}_h)\cdot \nabla\mathbf{w}_h\,\mathrm{d} \mathbf{x},
\\
I_2 &= -\sum_{e \in \mathcal{E}_h}\int_{e}\{\nabla S(\mathbf{z}_h)\} \cdot \mathbf{n}[[\mathbf{w}_h]]\,\mathrm{d}s,\\
I_3 &=
 \sum_{e \in \mathcal{E}_h}\int_{e}\{\nabla\mathbf{w}_h\} \cdot \mathbf{n}[[S(\mathbf{z}_h)]]\,\mathrm{d} s \\
 I_4 &= \mu_e \sum_{e \in \mathcal{E}_h} \int_{e}[[S(\mathbf{z}_h)]] \cdot [[\mathbf{w}_h]]\,\mathrm{d}s.
\end{aligned}
\end{equation}

For the gradient term, Cauchy--Schwarz inequality directly gives
\begin{equation}\notag
|I_1|
\le
\left(
\sum_{K\in Q_h}\|\nu_T\nabla (m_{1,h},m_{2,h})\|_{L^2(K)}^2
\right)^{1/2}
\left(
\sum_{K\in Q_h}\|\nabla (w_{1,h},w_{2,h})\|_{L^2(K)}^2
\right)^{1/2}.
\end{equation}
Since $\sigma>0$, this contribution is strictly bounded by $\nu_T\|\mathbf{z}_h\|_{\mathrm{DG}}\|\mathbf{w}_h\|_{\mathrm{DG}}$.

For the interface terms, applying Cauchy--Schwarz inequality on each face yields
\begin{equation}\notag
|I_2|
\le
\left(
\sum_{e\in \mathcal{E}_h}\frac{\nu_T^2h_e^\beta}{\sigma}\|\{\nabla (m_{1,h},m_{2,h})\}\cdot \mathbf{n}\|_{L^2(e)}^2
\right)^{1/2}
\left(
\sum_{e\in \mathcal{E}_h}\frac{\sigma}{h_e^\beta}\|[[ (w_{1,h},w_{2,h})]]\|_{L^2(e)}^2
\right)^{1/2}.
\end{equation}
Here, due to the fact that $h_e \le h_K \triangleq \sup\limits_{p, q \in K}\|p-q\| \le h$, and assuming without loss of generality that $\beta \ge 1$ and $h_K \le 1$, the first factor is controlled through the discrete trace inequality and local mesh regularity:
\begin{equation}\notag
\begin{aligned}
&\sum_{e\in \mathcal{E}_h}\frac{\nu_T^2h_e^\beta}{\sigma}\|\{\nabla (m_{1,h},m_{2,h})\}\cdot \mathbf{n}\|_{L^2(e)}^2\\
\le& \sum_{e \in \mathcal{E}_h} \frac{\nu_T^2h_e^\beta}{\sigma}\sum_{K \in \{K^+,K^-\}}\frac{1}{2} \| \left.\nabla (m_{1,h},m_{2,h})\cdot \mathbf{n} \right|_{K}\|_{L^2(e)}^2  \\
\le& \frac{\nu_T^2C_k^2}{2\sigma}\sum_{e \in \mathcal{E}_h} h_e^\beta (h_{K^+}^{-1} + h_{K^-}^{-1}) \sum_{K \in \{K^+,K^-\}} \|\nabla (m_{1,h},m_{2,h})\|_{L^2(K)}^2 \\
\le& 
\frac{\nu_T^2C_k^2 n_0}{\sigma} \sum_{K\in Q_h} \|\nabla (m_{1,h},m_{2,h})\|_{L^2(K)}^2,
\end{aligned}
\end{equation}
where $n_0$ denotes the maximum number of neighbors an element can have, and the constant $C_k$ derived from the trace inequality is independent of $h_K$ and $\mathbf{z}_h$, but depends on the polynomial degree $k$. Then, the term $|I_2|$ is bounded by $\frac{\nu_T C_k \sqrt{n_0}}{\sqrt{\sigma}}\|\mathbf{z}_h\|_{\mathrm{DG}}\|\mathbf{w}_h\|_{\mathrm{DG}}$.

A similar argument based on the trace inequality yields the same bound for $I_3$. Finally, evaluating the penalty term $I_4$ with $\mu_e = \sigma/h_e^\beta$ satisfies
\[
|I_4|
\le \nu_T
\left(
\sum_{e\in \mathcal{E}_h}\frac{\sigma}{h_e^\beta}
\|[[(m_{1,h},m_{2,h})]]\|_{L^2(e)}^2
\right)^{1/2}
\left(
\sum_{e \in \mathcal{E}_h}\frac{\sigma}{h_e^\beta}
\|[[(w_{1,h},w_{2,h})]]\|_{L^2(e)}^2
\right)^{1/2}.
\]
Collecting these bounds proves the continuity estimate \eqref{eq13}, with the combined constant $C_b=C \nu_T \cdot \max \left\{1,\frac{C_k \sqrt{n_0}}{\sqrt{\sigma}}\right\}$.

Setting $\mathbf w_h=\mathbf z_h$ makes the two nonsymmetric consistency terms cancel exactly, leaving the volume and penalty parts in \eqref{eq14},
therefore leading to the coercivity.
\end{proof}

\end{document}
